\chapter*{Streszczenie}
\thispagestyle{plain}

Dzisiejszy rynek podzespołów komputerowych rozwija się w bardzo dynamiczny sposób, co przyczynia się do coraz większej liczby ich możliwych typów, modeli oraz konfiguracji.
Zbudowanie komputera osobistego może więc wydawać się skomplikowanym zadaniem, biorąc pod uwagę kompatybilność sprzętu, a także różnice w budowie poszczególnych podzespołów.
Nauka montażu komputera na rzeczywistym sprzęcie niesie również ryzyko uszkodzenia drogich komponentów, w szczególności podczas montażu delikatnych elementów, takich jak procesor.
Tradycyjne metody nauki, takie jak książki i materiały wideo, nie zapewniają niezbędnej wiedzy oraz praktycznego doświadczenia.
Prowadzi to do rosnącego zapotrzebowania na interaktywne środowisko edukacyjne, pozwalające na naukę budowy komputera bez ponoszenia ryzyka finansowego.
Odpowiedzią na wymienione problemy może być aplikacja wykorzystująca technologię rzeczywistości wirtualnej, wprowadzająca użytkowników w tematykę montażu komputera.

Celem projektu inżynierskiego "Wirtualny montaż i wizualizacja działania komputera osobistego" było zaprojektowanie oraz implementacja wirtualnego trenażera operującego w jaskiniach rzeczywistości wirtualnej, znajdujących się w Laboratorium Zanurzonej Wizualizacji Przestrzennej Politechniki Gdańskiej.
W wyniku prac powstał symulator budowy komputera, oferujący użytkownikom wirtualny montaż licznych konfiguracji sprzętowych, który uwzględnia kompatybilność podzespołów oraz prowadzi ich przez pełen proces w realistyczny sposób.
Aplikacja zapewnia bezpieczne środowisko do nauki budowy komputera, w którym nie występuje ryzyko uszkodzenia drogich komponentów, a możliwość swobodnego powtarzania wykonanych czynności prowadzi do utrwalania nabytej wiedzy.
Wykorzystanie technologii rzeczywistości wirtualnej pozwoliło na dokładne wykrywanie oraz odwzorowanie pozycji rąk i głowy użytkownika, zapewniając mu wysoki poziom immersji w utworzonym środowisku.
W trenażerze zaimplementowane zostały niezbędne mechaniki pełnej kontroli nad podzespołami komputerowymi, pozwalające na ich swobodne przenoszenie oraz łączenie.
Aby zapewnić wsparcie użytkownikowi, utworzone zostały elementy interfejsu aplikacji, które prowadzą go przez każdy etap budowy wirtualnego komputera w przystępny sposób.
Symulator został wytworzony przy użyciu silnika gier Unreal Engine 5, który pozwolił na bezpośrednią współpracę z systemami śledzenia ruchu oraz synchronizacji wieloekranowej projekcji.

Symulator został dokładnie przetestowany w jaskiniach rzeczywistości wirtualnej, a także przeprowadzone zostały badania oceniające jego intuicyjność oraz wartość edukacyjną.
Testy z udziałem użytkowników wykazały ich wysoki poziom satysfakcji, potwierdzając skuteczność trenażera jako narzędzia edukacyjnego.
Dodatkowo przeprowadzono badania porównawcze symulatora z dostępnym na rynku alternatywnym rozwiązaniem, które wykazały zbliżony lub wyższy poziom satysfakcji użytkowników aplikacji rzeczywistości wirtualnej.
Uczestnicy badań pozytywnie ocenili intuicyjność trenażera oraz wykazali większą pewność siebie w zakresie samodzielnego montażu komputera.

Praca inżynierska była realizowana przez zespół czterech studentów.
Każdemu członkowi zespołu przypisane zostały następujące zadania podczas implementacji trenażera oraz realizacji pracy dyplomowej:
\begin{itemize}
    \item \textbf{Jakub Jędrzejczyk} - odpowiadał za zarządzanie zespołem projektowym oraz implementację mechanik kontrolowania podzespołów i montażu komputera oraz za stworzenie czytelnego interfejsu użytkownika. Jest autorem rozdziału \ref{chap:Testing} oraz podrozdziałów \ref{sec:GameEngineSelection}, \ref{chap:GUI}, \ref{chap:ObjectHolding}, \ref{chap:Dtrack}, \ref{chap:ConnectingComponents}, \ref{chap:PCValidation}, \ref{sec:FuturePlans}.
    \item \textbf{Jakub Kinder} - odpowiadał za zaprojektowanie i wytworzenie menedżera procesu budowy komputera, śledzącego jego postęp oraz walidującego utworzoną konfigurację. Jest autorem rozdziałów \ref{chap:Introduction} i \ref{chap:FieldIntroduction}, a także podrozdziałów \ref{chap:buildmanager}. 
    \item \textbf{Hubert Leszczyński} - odpowiadał za zaprojektowanie systemu, dekorację głównej sceny aplikacji, a także wytworzenie trójwymiarowych grafik. Zaimplementował również mechanizm dynamicznego tworzenia podzespołów. Jest autorem rozdziału \ref{chap:project} oraz podrozdziałów \ref{chap:models}, \ref{chap:spawners}, \ref{chap:objectConfig} i \ref{chap:sceneDesign}.
    \item \textbf{Mateusz Merda} - odpowiadał za wytworzenie elementów posiadających zaawansowane konfiguracje fizyki i kolizji. Jest autorem rozdziału 3, a także podrozdziałów \ref{subsec:twohands}, \ref{sec:cableconnector}, \ref{sec:results}, \ref{sec:conclusions}.
\end{itemize}

W przyszłości projekt może zostać rozwinięty o dokładne mechanizmy wizualizacji działania komputera osobistego, a także o większą liczbę dostępnych podzespołów i konfiguracji sprzętowych.
Słowa kluczowe: Rzeczywistość wirtualna, Wytwarzanie aplikacji, Wizualizacja, Budowa komputera
Dziedzina nauki i techniki, zgodnie z wymogami OECD: 1.2 Nauki o komputerach i informatyka